\section{Experimental Setup and Results}
\subsection{Word-level Tokenization Results}

\subsubsection{Experiment 1: Baseline Vocabulary (5,000 Tokens)}
We standardized hyperparameters (Chunk Size: 200, Seed: 42, Subset: 1.0) and evaluated a constrained vocabulary of 5,000 tokens.

\begin{table}[H]
\centering
\caption{Detailed Tokenization Results (Vocab Size = 5,000)}
\renewcommand{\arraystretch}{1.1}
\begin{tabular}{|l|l|c|c|c|c|}
\hline
\textbf{Split} & \textbf{Metric} & \textbf{Enwik8} & \textbf{One Billion} & \textbf{Text8} & \textbf{WikiText-103} \\ \hline
\multirow{3}{*}{Train} 
 & Avg Seq Length & 194.353 & 26.652 & 199.997 & 86.713 \\ \cline{2-6}
 & OOV Rate (\%) & 0 & 0 & 0 & 0 \\ \cline{2-6}
 & Compression Ratio & 0.2098 & 0.1965 & 0.1700 & 0.1910 \\ \hline
\multirow{3}{*}{Test}  
 & Avg Seq Length & 193.990 & 26.671 & 200.000 & 82.898 \\ \cline{2-6}
 & OOV Rate (\%) & 15.555 & 11.718 & 16.181 & 14.844 \\ \cline{2-6}
 & Compression Ratio & 0.2095 & 0.1965 & 0.1698 & 0.1898 \\ \hline
\end{tabular}
\end{table}

\textbf{Insight:} A 5,000-word limit causes severe semantic loss. The test OOV rate aggressively spikes to $\sim$16\% on Text8 and $\sim$15.5\% on Enwik8, rendering this configuration unviable for language modeling.

% ------------------------------------------------------
% EXPERIMENT 2: VOCAB 10000
% ------------------------------------------------------
\subsubsection{Experiment 2: Expanding Vocabulary to 10,000 Tokens}

\begin{table}[H]
\centering
\caption{Detailed Tokenization Results (Vocab Size = 10,000)}
\renewcommand{\arraystretch}{1.1}
\begin{tabular}{|l|l|c|c|c|c|}
\hline
\textbf{Split} & \textbf{Metric} & \textbf{Enwik8} & \textbf{One Billion} & \textbf{Text8} & \textbf{WikiText-103} \\ \hline
\multirow{3}{*}{Train} 
 & Avg Seq Length & 194.353 & 26.652 & 199.997 & 86.713 \\ \cline{2-6}
 & OOV Rate (\%) & 0 & 0 & 0 & 0 \\ \cline{2-6}
 & Compression Ratio & 0.2098 & 0.1965 & 0.1700 & 0.1910 \\ \hline
\multirow{3}{*}{Test}  
 & Avg Seq Length & 193.990 & 26.671 & 200.000 & 82.898 \\ \cline{2-6}
 & OOV Rate (\%) & 10.432 & 6.853 & 10.375 & 9.658 \\ \cline{2-6}
 & Compression Ratio & 0.2095 & 0.1965 & 0.1698 & 0.1898 \\ \hline
\end{tabular}
\end{table}

\textbf{Insight:} Doubling the vocabulary halves the OOV rate for the One Billion Word dataset (6.85\%). However, a $\sim$10\% OOV rate on the other datasets still leaves a significant portion of text uninterpretable.

% ------------------------------------------------------
% EXPERIMENT 3: VOCAB 50000
% ------------------------------------------------------
\subsubsection{Experiment 3: Scaling Vocabulary to 50,000 Tokens}

\begin{table}[H]
\centering
\caption{Detailed Tokenization Results (Vocab Size = 50,000)}
\renewcommand{\arraystretch}{1.1}
\begin{tabular}{|l|l|c|c|c|c|}
\hline
\textbf{Split} & \textbf{Metric} & \textbf{Enwik8} & \textbf{One Billion} & \textbf{Text8} & \textbf{WikiText-103} \\ \hline
\multirow{3}{*}{Train} 
 & Avg Seq Length & 194.353 & 26.652 & 199.997 & 86.713 \\ \cline{2-6}
 & OOV Rate (\%) & 0 & 0 & 0 & 0 \\ \cline{2-6}
 & Compression Ratio & 0.2098 & 0.1965 & 0.1700 & 0.1910 \\ \hline
\multirow{3}{*}{Test}  
 & Avg Seq Length & 193.990 & 26.671 & 200.000 & 82.898 \\ \cline{2-6}
 & OOV Rate (\%) & 3.408 & 1.314 & 2.773 & 2.765 \\ \cline{2-6}
 & Compression Ratio & 0.2095 & 0.1965 & 0.1698 & 0.1898 \\ \hline
\end{tabular}
\end{table}

\textbf{Insight:} At 50,000 tokens, the tokenizer achieves acceptable semantic retention (OOV $\le$ 3.4\%) while maintaining high compression efficiency (0.16–0.20), resolving the missing vocabulary bottleneck.

% ======================================================
% CHARACTER-LEVEL TOKENIZATION EXPERIMENTS
% ======================================================
\subsection{Character-level Tokenization Results}
This approach treats every individual character as a distinct token, fundamentally altering vocabulary constraints.

\begin{table}[H]
\centering
\caption{Character-level Tokenization Results across Datasets}
\renewcommand{\arraystretch}{1.1}
\begin{tabular}{|l|l|c|c|c|c|}
\hline
\textbf{Split} & \textbf{Metric} & \textbf{Enwik8} & \textbf{One Billion} & \textbf{Text8} & \textbf{WikiText-103} \\ \hline
\multicolumn{2}{|c|}{\textbf{Absolute Vocab Size}} & \textbf{52} & \textbf{52} & \textbf{31} & \textbf{52} \\ \hline\hline
\multirow{3}{*}{Test}  
 & Avg Seq Length & 925.723 & 135.270 & 1178.084 & 431.722 \\ \cline{2-6}
 & OOV Rate (\%) & 0 & 0 & 0 & 0 \\ \cline{2-6}
 & Compression Ratio & 1.000 & 1.000 & 1.000 & 1.000 \\ \hline
\end{tabular}
\end{table}

\textbf{Insights:}
\begin{itemize}
    \item \textbf{Microscopic Vocabulary:} Uses only 31-52 tokens, drastically reducing the embedding memory footprint.
    \item \textbf{0\% OOV \& Zero Compression:} Mathematically guarantees no missing words, but offers no text compression (Ratio = 1.0).
    \item \textbf{Sequence Inflation:} Sequence lengths explode (e.g., $\sim$1178 on Text8). This places an immense computational burden on recurrent models (LSTMs) to capture long-term context.
\end{itemize}

% ======================================================
% BPE TOKENIZATION EXPERIMENTS
% ======================================================
\subsection{Byte Pair Encoding (BPE) Results}
BPE dynamically merges frequent character pairs into subwords. We evaluated it across 16k, 32k, and 50k vocabulary sizes (Note: One Billion Word saturated at 38,379).

\begin{table}[H]
\centering
\caption{BPE Test Metrics: Enwik8 \& One Billion Word}
\renewcommand{\arraystretch}{1.1}
\begin{tabular}{|l|l|c|c|c|}
\hline
\textbf{Dataset} & \textbf{Metric (Test Split)} & \textbf{Vocab=16k} & \textbf{Vocab=32k} & \textbf{Vocab=50k} \\ \hline
\multirow{2}{*}{Enwik8} 
 & Avg Seq Length & 213.330 & 201.975 & 197.240 \\ \cline{2-5}
 & Compression Ratio & 0.2304 & 0.2182 & 0.2131 \\ \hline
\multirow{2}{*}{One Billion*} 
 & Avg Seq Length & 28.984 & 28.035 & 27.829 \\ \cline{2-5}
 & Compression Ratio & 0.2143 & 0.2073 & 0.2057 \\ \hline
\multicolumn{5}{l}{\footnotesize \textit{*One Billion Word vocab saturated at 38,379 instead of 50,000.}}
\end{tabular}
\end{table}

\begin{table}[H]
\centering
\caption{BPE Test Metrics: Text8 \& WikiText-103}
\renewcommand{\arraystretch}{1.1}
\begin{tabular}{|l|l|c|c|c|}
\hline
\textbf{Dataset} & \textbf{Metric (Test Split)} & \textbf{Vocab=16k} & \textbf{Vocab=32k} & \textbf{Vocab=50k} \\ \hline
\multirow{2}{*}{Text8} 
 & Avg Seq Length & 226.540 & 213.795 & 208.952 \\ \cline{2-5}
 & Compression Ratio & 0.1923 & 0.1815 & 0.1774 \\ \hline
\multirow{2}{*}{WikiText-103} 
 & Avg Seq Length & 92.654 & 88.525 & 86.819 \\ \cline{2-5}
 & Compression Ratio & 0.2146 & 0.2051 & 0.2011 \\ \hline
\end{tabular}
\end{table}

\textbf{Unified BPE Insights:}
\begin{itemize}
    \item \textbf{Complete OOV Elimination:} Across all datasets, the OOV rate is strictly \textbf{0\%}. BPE successfully resolves the vocabulary bottleneck of Word-level tokenization by breaking rare words into known subwords.
    \item \textbf{Optimal Compression:} At a 50k vocabulary, BPE achieves compression ratios (0.17 - 0.21) nearly identical to pure Word-level methods. It compresses text highly efficiently while preventing sequence inflation (unlike Character-level).
    \item \textbf{Diminishing Returns:} Expanding the vocabulary from 16k to 32k yields a significant reduction in sequence length. However, scaling further to 50k provides only marginal improvements, indicating that an optimal balance lies around 30,000 tokens.
\end{itemize}

\textbf{Conclusion:} BPE serves as the optimal bridge. It preserves the memory efficiency of character-level models (no OOV) while retaining the computational speed and short sequence lengths of word-level models.